% Options for packages loaded elsewhere
\PassOptionsToPackage{unicode}{hyperref}
\PassOptionsToPackage{hyphens}{url}
\documentclass[
  10pt,
]{article}
\usepackage{xcolor}
\usepackage[margin=25mm]{geometry}
\usepackage{amsmath,amssymb}
\setcounter{secnumdepth}{-\maxdimen} % remove section numbering
\usepackage{iftex}
\ifPDFTeX
  \usepackage[T1]{fontenc}
  \usepackage[utf8]{inputenc}
  \usepackage{textcomp} % provide euro and other symbols
\else % if luatex or xetex
  \usepackage{unicode-math} % this also loads fontspec
  \defaultfontfeatures{Scale=MatchLowercase}
  \defaultfontfeatures[\rmfamily]{Ligatures=TeX,Scale=1}
\fi
\usepackage{lmodern}
\ifPDFTeX\else
  % xetex/luatex font selection
\fi
% Use upquote if available, for straight quotes in verbatim environments
\IfFileExists{upquote.sty}{\usepackage{upquote}}{}
\IfFileExists{microtype.sty}{% use microtype if available
  \usepackage[]{microtype}
  \UseMicrotypeSet[protrusion]{basicmath} % disable protrusion for tt fonts
}{}
\makeatletter
\@ifundefined{KOMAClassName}{% if non-KOMA class
  \IfFileExists{parskip.sty}{%
    \usepackage{parskip}
  }{% else
    \setlength{\parindent}{0pt}
    \setlength{\parskip}{6pt plus 2pt minus 1pt}}
}{% if KOMA class
  \KOMAoptions{parskip=half}}
\makeatother
\setlength{\emergencystretch}{3em} % prevent overfull lines
\providecommand{\tightlist}{%
  \setlength{\itemsep}{0pt}\setlength{\parskip}{0pt}}
\usepackage{bookmark}
\IfFileExists{xurl.sty}{\usepackage{xurl}}{} % add URL line breaks if available
\urlstyle{same}
\hypersetup{
  pdftitle={RA-11: probability-sensitive product probes and diagonal correction},
  pdfauthor={Sidney Holden},
  hidelinks,
  pdfcreator={LaTeX via pandoc}}

\title{RA-11: probability-sensitive product probes and diagonal
correction}
\author{Sidney Holden}
\date{14 September 2026}

\usepackage{fvextra,xurl}
\DefineVerbatimEnvironment{verbatim}{Verbatim}{breaklines=true,fontsize=\small}
\begin{document}
\maketitle

\textbf{Affiliation:} Center for Computational Biology, Flatiron
Institute, Simons Foundation.

\textbf{Submission disposition:} Partially resolved. Theorems 2.1 and
2.2 and Corollary 2.3 pass for the new upper bounds and
estimator-specific rate. The unrestricted joint minimax characterization
remains open.

\href{https://github.com/ajt60gaibb/OpenProblemsInNLA/blob/main/reviews/2026-09-14-further-submissions/RA-11-review.md}{Independent
informal AI-agent review}. ChatGPT assistance is disclosed. No Lean
verification, external human peer review or formal verification is
claimed. This dated notice supersedes historical statements about
pending review; inherited claims are accepted only within the linked
review scope.

\textbf{Research continuation --- round 4. Status: partial, not a full
solution.}

\subsection{Abstract}\label{abstract}

We consider relative trace estimation of an arbitrary real
positive-semidefinite matrix of order N=n\^{}q, when an oracle returns
the full matrix-vector response only on real product inputs. Two new
upper bounds are established. A fractional-moment analysis of
product-spherical empirical averaging improves the fixed-accuracy
exponential rate from log(2n/(n+1)) to D\_n=log n+1-H\_n.~The rate D\_n
is sharp for that empirical-average estimator, but not a minimax lower
bound for the full oracle. Separately, a two-stage diagonal-control
estimator has variance at most
2N\textbar\textbar M\textbar\textbar\_F\^{}2/(ms) after m training
queries and s coordinate corrections, giving O(sqrt(N)/epsilon) queries
with finite-bit, nonadaptive real inputs. Both results hold for every
PSD input. The remaining unrestricted lower/upper gap is stated
explicitly.

\subsection{1. Model and relationship to the previous
packages}\label{model-and-relationship-to-the-previous-packages}

Fix integers n,q \textgreater= 2 and 0 \textless{} epsilon \textless{}
1/2. The unknown real symmetric matrix M has order N=n\^{}q and
satisfies M \textgreater= 0. One oracle call supplies real vectors
v\_1,\ldots,v\_q in R\^{}n and returns the entire vector

\begin{verbatim}
M(v_1 tensor ... tensor v_q).
\end{verbatim}

The matrix M itself need not factor. Define Q(n,q,epsilon) as the
minimum worst-case number of calls needed by a randomized algorithm
that, for every fixed M, returns T with

\begin{verbatim}
P(|T - tr M| <= epsilon tr M) >= 2/3.
\end{verbatim}

Finite exact computation is uncharged; no query-conditioning bound is
imposed. These are the model requirements used in the preceding
manuscripts and in the repository's RA-11 statement {[}R1{]}. All new
proofs below concern this stronger full-vector oracle, not a conditioned
scalar oracle.

Write tau=tr M. When tau=0, PSD implies M=0 and the estimators below
return zero exactly.

The previous packages {[}R2--R4{]} report, among other results, the
retained lower bound

\begin{verbatim}
Q >= (1/80000) max_{1<=b<=q} min{n^b, sqrt(floor(q/b))/epsilon}.
\end{verbatim}

They also give fixed-q optimality and several upper bounds. Those
results are inherited, not newly proved here. This manuscript supplies
proofs of its own new upper bounds from elementary calculations and
explicitly distinguishes them from the inherited assertions.

\subsection{2. Main results}\label{main-results}

For 1 \textless{} p \textless= 2 define

\begin{verbatim}
C(n,p) = n^p Gamma(p+1) Gamma(n) / Gamma(n+p).
\end{verbatim}

Equivalently, with t=p-1,

\begin{verbatim}
C(n,1+t) = n^t product_{k=2}^n k/(k+t).
\end{verbatim}

\subsubsection{Theorem 2.1 --- fractional-moment upper
bound}\label{theorem-2.1-fractional-moment-upper-bound}

For any 1 \textless{} p \textless= 2, an ideal complex-spherical
estimator can be implemented with real product queries using

\begin{verbatim}
(q+1) ceil[(6 C(n,p)^q / epsilon^p)^(1/(p-1))]
\end{verbatim}

calls, and its success probability is at least 2/3 for every real PSD M.
A version with a deterministically finite number of fair random bits
exists with no more than twice the number of samples in this display. In
either case the query list can be fixed before responses are observed.
The exact basis bound N may always be used instead.

The ideal formulation uses continuous random probes, as is customary in
exact-real oracle analysis. Section 6 gives the finite-bit reduction; it
does not assume floating-point Gaussian simulation is exact.

\subsubsection{Theorem 2.2 --- diagonal-control upper
bound}\label{theorem-2.2-diagonal-control-upper-bound}

There is an explicit bounded-fair-bit, nonadaptive algorithm with

\begin{verbatim}
Q(n,q,epsilon) <= min{N, 2 ceil[sqrt(6N)/epsilon]}.
\end{verbatim}

When N is a power of two, the constant 6 can be replaced by 3. Only real
Rademacher-product inputs and real coordinate-product inputs are used.
This construction does not use interpolation, Gaussian sampling, or a
low-rank approximation subroutine.

\subsubsection{Corollary 2.3 --- fixed-accuracy rate of the new first
branch}\label{corollary-2.3-fixed-accuracy-rate-of-the-new-first-branch}

Put

\begin{verbatim}
D_n = log n + 1 - H_n,
V_n = sum_{k=2}^n 1/k^2,
A = log(6/epsilon).
\end{verbatim}

If q V\_n \textgreater= 2A, the ideal sample count in Theorem 2.1 is at
most

\begin{verbatim}
ceil[epsilon^(-1) exp(q D_n + sqrt(2 q V_n A))].
\end{verbatim}

The finite-bit version uses at most twice this sample count. Each sample
costs q+1 real queries. In particular, at n=2 and epsilon=1/4, for
q\textgreater=26,

\begin{verbatim}
Q(2,q,1/4)
  <= min{2^q,
         2(q+1) ceil[4 exp((log 2 - 1/2)q
                               + sqrt((q/2)log 24))]}.
\end{verbatim}

Thus the fixed-accuracy exponential base of this upper branch is
2/sqrt(e), approximately 1.21306. This is strictly below 4/3, the base
supplied by the preceding complex-spherical variance calculation. It is
still an exponential upper bound, and it does not match the known
unrestricted lower scale.

\subsection{3. Convex-order reduction to a scalar product
distribution}\label{convex-order-reduction-to-a-scalar-product-distribution}

Let z\_j be independent and uniform on the complex sphere of radius
sqrt(n), and set z=z\_1 tensor \ldots{} tensor z\_q. Let W=nB, where B
has the Beta(1,n-1) distribution. Then EW=1. Define
P\_q=product\_\{j=1\}\^{}q W\_j for independent copies W\_j.

Recall that X is dominated by Y in convex order when E phi(X)\textless=E
phi(Y) for every convex function for which the expectations exist, and
EX=EY.

\subsubsection{Lemma 3.1 --- one local
factor}\label{lemma-3.1-one-local-factor}

For any unit u in C\^{}n, \textbar u* z\_1\textbar\^{}2 has law W. If
A\textgreater=0 and tr A=1, then z\_1* A z\_1 is dominated by W in
convex order.

Proof. Unitary invariance gives the first assertion. Diagonalize
A=sum\_r lambda\_r u\_r u\_r*, where lambda\_r\textgreater=0 and sum
lambda\_r=1. Pointwise Jensen gives

\begin{verbatim}
phi(sum_r lambda_r |u_r* z_1|^2)
   <= sum_r lambda_r phi(|u_r* z_1|^2).
\end{verbatim}

Take expectations. Every summand on the right has the W distribution.
Isotropy supplies equality of means. QED.

\subsubsection{Theorem 3.2 --- all-PSD convex-order
domination}\label{theorem-3.2-all-psd-convex-order-domination}

For every Hermitian PSD M with positive trace tau,

\begin{verbatim}
(z* M z)/tau  is dominated by  P_q in convex order.
\end{verbatim}

Proof. Induct on q, including arbitrary PSD matrices in the induction
statement. For a rank-one matrix uu* with
\textbar\textbar u\textbar\textbar=1, condition on the first q-1
factors. The contraction against these factors is a vector a in
C\^{}n.~The conditional quadratic response has law
\textbar\textbar a\textbar\textbar\^{}2 W, with W independent of the
conditioning variables: this follows from local unitary invariance,
including the case a=0. The quantity
\textbar\textbar a\textbar\textbar\^{}2 is a product-probe quadratic
form of the PSD partial trace of uu*, on q-1 modes, with trace one. By
induction it is convex-order dominated by P\_\{q-1\}. Multiplication by
an independent nonnegative W preserves this order, since x
-\textgreater{} E phi(Wx) is convex. This proves the assertion for rank
one. For general trace-one M, use its spectral decomposition and the
same pointwise Jensen argument as in Lemma 3.1. Rescaling handles tau.
QED.

Two consequences should not be confused. First,

\begin{verbatim}
E(z* M z)^p <= tau^p C(n,p)^q
\end{verbatim}

for p\textgreater=1. Second, replacing independent copies one at a time
shows that the empirical average of normalized responses is convex-order
dominated by the empirical average of independent P\_q variables. This
gives bounds for convex losses, but it does not say that every tail
probability is maximized by a rank-one input.

The moment formula follows directly from the beta integral:

\begin{verbatim}
EW^p = n^p (n-1) integral_0^1 b^p(1-b)^(n-2) db
      = C(n,p).
\end{verbatim}

No assumption about a tensor-product factorization of M was made.

\subsection{4. Bounded-probability analysis and real-query
accounting}\label{bounded-probability-analysis-and-real-query-accounting}

\subsubsection{Lemma 4.1 --- an elementary mean
bound}\label{lemma-4.1-an-elementary-mean-bound}

Let X\_1,\ldots,X\_m be independent nonnegative identically distributed
variables with mean mu and finite p-th moment, 1\textless p\textless=2.
Then

\begin{verbatim}
E |m^(-1) sum_i X_i - mu|^p <= 2 E X_1^p / m^(p-1).
\end{verbatim}

Proof. Introduce independent copies X\_i'. Jensen bounds the centered
sum's p-th moment by that of sum\_i(X\_i-X\_i'). The independent
differences are symmetric, so independent random signs may be inserted
without changing their joint distribution. Conditional on the
differences d\_i, the sign average is bounded by

\begin{verbatim}
E_sign |sum_i sign_i d_i|^p
   <= (sum_i d_i^2)^(p/2)
   <= sum_i |d_i|^p.
\end{verbatim}

The first inequality is monotonicity of moments from p to 2; the second
uses p\textless=2. For nonnegative a,b,
\textbar a-b\textbar{}\textsuperscript{p\textless=a}p+b\^{}p.~Taking
expectations and dividing by m\^{}p proves the claim. QED.

Apply this lemma with X=z* M z. Theorem 3.2 and Markov's inequality give

\begin{verbatim}
P(|mean X - tau| > epsilon tau)
   <= 2 C(n,p)^q / (m^(p-1) epsilon^p).
\end{verbatim}

The sample count in Theorem 2.1 makes the right side at most 1/3.

\subsubsection{Lemma 4.2 --- real simulation of a complex product
response}\label{lemma-4.2-real-simulation-of-a-complex-product-response}

For arbitrary complex factors g\_j=a\_j+i b\_j, consider the vector
polynomial

\begin{verbatim}
v(t) = tensor_{j=1}^q (a_j+t b_j).
\end{verbatim}

Its degree is at most q. Query Mv(t) at q+1 distinct real values, for
example 0,1,\ldots,q, and interpolate its value at t=i. This returns
M(g\_1 tensor \ldots{} tensor g\_q) using only real product inputs.

Explicitly, the weight for node j is

\begin{verbatim}
ell_j(i)= product_{0<=h<=q, h!=j} (i-h)/(j-h).
\end{verbatim}

These weights are Gaussian rational. There are finitely many exact
arithmetic operations, although floating-point interpolation can be
badly conditioned.

To sample a spherical probe without taking a square root in unknown-data
arithmetic, draw nonzero complex Gaussian factors g\_j and query their
raw product. If g= tensor\_j g\_j, return

\begin{verbatim}
[product_j n/(g_j* g_j)] (g* M g).
\end{verbatim}

This equals z* M z for the associated normalized spherical factors. All
normalizing weights are real scalar divisions. The same raw-query
implementation applies to the finite-support construction in Section 6.

The query list is nonadaptive: all factors and interpolation nodes are
selected before any oracle response.

\subsection{5. The exponential rate and its exact
scope}\label{the-exponential-rate-and-its-exact-scope}

Let t=p-1. A useful exact identity is

\begin{verbatim}
f_n(t) := log C(n,1+t)
         = t log n - sum_{k=2}^n log(1+t/k).
\end{verbatim}

Hence

\begin{verbatim}
f_n'(0)=D_n,
f_n''(t)=sum_{k=2}^n 1/(k+t)^2 <= V_n.
\end{verbatim}

Taylor's theorem with an integral remainder gives

\begin{verbatim}
f_n(t) <= D_n t + (V_n/2)t^2,  0<=t<=1.
\end{verbatim}

The logarithm of the unrounded ideal sample count is at most

\begin{verbatim}
q D_n + log(1/epsilon) + A/t + (q V_n/2)t,
A=log(6/epsilon).
\end{verbatim}

When sqrt(2A/(qV\_n))\textless=1, choose this value of t to obtain
Corollary 2.3. Otherwise choose t=1; the original formula of Theorem 2.1
remains valid at every parameter value.

\subsubsection{Theorem 5.1 --- sharp rate for empirical averaging, not
for Q}\label{theorem-5.1-sharp-rate-for-empirical-averaging-not-for-q}

Fix n and 0\textless epsilon\textless1. Consider ideal spherical
empirical averaging with m\_q independent probes. If

\begin{verbatim}
limsup_{q->infinity} (log m_q)/q < D_n,
\end{verbatim}

then on a trace-one rank-one product matrix its empirical mean converges
to zero in probability. Conversely, for any fixed r\textgreater D\_n,
taking m\_q=ceil(exp(rq)) makes the empirical mean converge to one in
probability uniformly over all trace-one PSD matrices of order n\^{}q.

Proof of the converse. Since f\_n(t)/t tends to D\_n as t decreases to
zero, choose a fixed t\textgreater0 with f\_n(t)/t\textless r. The bound
in Section 4 then tends to zero for each fixed epsilon, uniformly over
M.

Proof of the lower statement. For a rank-one product matrix, a
normalized response has exactly the P\_q distribution. Tilt the local
law of W by the density W, which is a probability density because EW=1.
Under this tilted law,

\begin{verbatim}
E_tilt log W = E[W log W] = D_n.
\end{verbatim}

Its log moment is integrable. For any a\textless D\_n, the law of large
numbers therefore gives

\begin{verbatim}
E[P_q 1{log P_q <= aq}]
   = P_tilt(sum_j log W_j <= aq) -> 0.
\end{verbatim}

Choose a strictly between the assumed upper exponential rate of m\_q and
D\_n.~A union bound and Markov's inequality give

\begin{verbatim}
P(any sampled P_q > exp(aq)) <= m_q exp(-aq) -> 0.
\end{verbatim}

On the complementary event, the empirical mean equals the average of the
truncated samples. That truncated average has expectation tending to
zero by the preceding tilted identity, and thus tends to zero in
probability. QED.

This theorem is an estimator-specific obstruction. It is emphatically
NOT an exponential lower bound for RA-11. If M=uu\^{}T and a query x
satisfies u\^{}T x != 0, then with y=Mx,

\begin{verbatim}
M = yy^T/(x^T y),
tr M = ||y||^2/(x^T y).
\end{verbatim}

A generic continuous product x has nonzero overlap almost surely with
any fixed nonzero u. Thus one full-vector query recovers the very
rank-one example used in the empirical-average lower bound. Confusing
these two oracle statements would invalidate a claimed complete
solution.

For n=2 one may also write W=2U with U uniform on {[}0,1{]}. Then

\begin{verbatim}
log P_q = q log 2 - G_q,
\end{verbatim}

where G\_q is Gamma with shape q and rate 1. Under the size-biased law
its rate is 2. These formulas allow scalar experiments at moderately
large q without constructing a 2\textsuperscript{q-by-2}q matrix. Such
experiments concern the empirical estimator only.

\subsection{6. Bounded-fair-bit reduction for the fractional-moment
bound}\label{bounded-fair-bit-reduction-for-the-fractional-moment-bound}

This section establishes existence of a finite exact implementation with
at most a factor two in sample count. It is separate from the
floating-point Gaussian diagnostic code.

\subsubsection{Lemma 6.1 --- finite isotropic angular
approximation}\label{lemma-6.1-finite-isotropic-angular-approximation}

For fixed n, p in (1,2{]}, and delta\textgreater0, there is a finite
dyadic-probability distribution on nonzero Gaussian-rational raw vectors
g in C\^{}n such that the formal normalized vector
z=sqrt(n)g/\textbar\textbar g\textbar\textbar{} satisfies

\begin{verbatim}
E zz* = I,
||z||^2 = n,
sup_{||u||=1} E|u* z|^(2p) <= C(n,p)+delta.
\end{verbatim}

It can be sampled using a deterministically bounded number of fair bits.
Raw query inputs and normalization squared weights require no square
roots.

Proof. Approximate a real standard Gaussian by a symmetric finite
distribution on nonzero rational half-grid points, with dyadic
probabilities. Let its support grow and mesh and probability error
decrease. Take 2n independent coordinates with that distribution,
grouping them into n complex coordinates. Every resulting raw vector is
nonzero. The distributions converge weakly to a nondegenerate complex
Gaussian. After normalization, their angular distributions converge
weakly to complex Haar measure, because normalization is continuous away
from zero and the limiting Gaussian has no atom at zero.

The finite laws are invariant under complex-coordinate permutations and
multiplication of any coordinate by i: real and imaginary coordinates
have the same symmetric law. These symmetries force the covariance of z
to be a scalar identity, and its deterministic squared norm n makes the
scalar one. This is exact, not approximate isotropy.

On the compact unit sphere, the functions z -\textgreater{} \textbar u*
z\textbar\^{}(2p), indexed by unit u, are uniformly bounded and
equicontinuous. A finite net in u and weak convergence imply uniform
convergence of their expectations. A sufficiently fine law therefore
gives the required moment bound. The finite laws can be chosen
effectively: Gaussian tail bounds, rational mesh bounds away from the
origin, and rational approximations to the one-dimensional bin
probabilities give explicit finite error certificates. Symmetric dyadic
masses are used in pairs, and the remaining mass is assigned
symmetrically. Each coordinate draw then uses a fixed number of bits.
QED.

For completeness, an elementary modulus suffices for the constructive
assertion. Outside an event where
\textbar\textbar G\textbar\textbar\textless a or a coordinate exceeds a
chosen truncation radius R, a coordinate rounding error at most h
changes the normalized direction by at most 2sqrt(2n)h/a. On the unit
sphere the directional moment function has Lipschitz constant at most 2p
n\^{}p.~The exceptional probability is bounded by
(sqrt(2/pi)a)\^{}(2n)+4n exp(-R\^{}2/2), plus the total variation error
in the finite scalar probabilities. Choose a small rational a, an
integer R, a rational h, and dyadic probability precision in that order.
These bounds tend to zero and can be compared to any positive rational
tolerance. This is finite preparatory computation; it need not be
efficient.

\subsubsection{Lemma 6.2 --- tensorizing a local moment
bound}\label{lemma-6.2-tensorizing-a-local-moment-bound}

If local probes are isotropic and satisfy E\textbar u*
z\_j\textbar{}\textsuperscript{(2p)\textless=h\textbar\textbar u\textbar\textbar{}}(2p),
then for q independent factors and every PSD M,

\begin{verbatim}
E(z* M z)^p <= h^q (tr M)^p.
\end{verbatim}

Proof. For a vector u with slices u\_i, condition on one factor and
apply its moment bound. Minkowski's inequality in L\_p gives

\begin{verbatim}
[E(sum_i |u_i* z_remaining|^2)^p]^(1/p)
   <= sum_i [E|u_i* z_remaining|^(2p)]^(1/p).
\end{verbatim}

Induction proves the vector bound
h\textsuperscript{q\textbar\textbar u\textbar\textbar{}}(2p); a spectral
decomposition and Minkowski give the PSD bound. QED.

Let t=p-1 and choose the finite law so that its local moment h obeys

\begin{verbatim}
h <= C(n,p) exp[t log(2)/q].
\end{verbatim}

Then h\^{}q \textless= 2\^{}t C(n,p)\^{}q. Doubling the ideal number of
samples compensates for this factor in the denominator m\^{}t of the
probability bound. Isotropy keeps the estimator exactly unbiased. This
proves the finite-bit part of Theorem 2.1. The complex factors are
implemented by Lemma 4.2 using raw Gaussian-rational vectors.

The finite law may be very large. This result makes no claim of
practical bit complexity, and the package does not include a certified
generator for that law.

\subsection{7. A two-stage diagonal
estimator}\label{a-two-stage-diagonal-estimator}

The following proof is independent of the spherical analysis.

\subsubsection{7.1 Training from full-vector
responses}\label{training-from-full-vector-responses}

For each training query draw independent local Rademacher vectors
r\_1,\ldots,r\_q in \{-1,+1\}\^{}n and set x=r\_1 tensor \ldots{} tensor
r\_q. Then x\_i\^{}2=1 and E x\_i x\_j=delta\_ij. Query y=Mx and form
the entire vector

\begin{verbatim}
a(x)_i = x_i y_i.
\end{verbatim}

Let d\_i=M\_ii and let a be the average of m independent such vectors.

\subsubsection{Lemma 7.1 --- simultaneous diagonal moment
identity}\label{lemma-7.1-simultaneous-diagonal-moment-identity}

\begin{verbatim}
E a = d,
E ||a-d||_2^2 = [||M||_F^2 - sum_i M_ii^2]/m.
\end{verbatim}

Proof. For one sample,

\begin{verbatim}
a(x)_i-d_i = sum_{j!=i} M_ij x_i x_j.
\end{verbatim}

Isotropy proves the mean identity. Since x\_i\^{}2=1, the squared error
has expectation sum\_\{j!=i\} M\_ij\^{}2; no fourth-moment independence
assumption is needed. Sum over i and use independence of training
samples. QED.

The full-vector response is crucial: one call supplies all N training
estimates. This does not charge N separate scalar queries.

\subsubsection{7.2 Finite-bit coordinate
corrections}\label{finite-bit-coordinate-corrections}

Choose a probability distribution p\_i\textgreater0 on the N coordinate
indices, independently of the training stage. Draw s independent indices
I\_l. Each e\_\{I\_l\} is itself a product of local coordinate vectors.
Query M e\_\{I\_l\} to obtain d\_\{I\_l\}, and return

\begin{verbatim}
T = sum_i a_i + (1/s) sum_{l=1}^s (d_{I_l}-a_{I_l})/p_{I_l}.
\end{verbatim}

Conditional on a, this estimator has mean tau. Its conditional variance
is

\begin{verbatim}
(1/s) [sum_i (d_i-a_i)^2/p_i - (sum_i(d_i-a_i))^2]
  <= (max_i 1/p_i)/s * ||d-a||_2^2.
\end{verbatim}

Since the conditional mean is exactly tau, there is no extra variance
term from the random training sum. Lemma 7.1 proves

\begin{verbatim}
Var T <= (max_i 1/p_i)/(ms) * ||M||_F^2.
\end{verbatim}

PSD gives \textbar\textbar M\textbar\textbar\_F\textless=tr M, so the
relative variance is bounded by (max\_i 1/p\_i)/(ms).

Uniform sampling gives max\_i 1/p\_i=N. For a bounded-bit rule valid for
every N, let B=ceil(log\_2 N), draw U uniformly from
\{0,\ldots,2\^{}B-1\}, and set I=U mod N. Its exact probability is

\begin{verbatim}
p_i = [floor(2^B/N) + 1{i < (2^B mod N)}]/2^B,
i=0,...,N-1.
\end{verbatim}

It satisfies p\_i\textgreater=1/(2N), uses exactly B bits, and has a
known rational probability for the correction weight. When N is a power
of two it is exactly uniform.

Taking m=s=ceil(sqrt(6N)/epsilon) gives relative variance at most
epsilon\^{}2/3. Chebyshev proves success probability at least 2/3. For a
power-of-two N take m=s=ceil(sqrt(3N)/epsilon). Use the cheaper exact
basis sweep whenever its N calls are fewer. This proves Theorem 2.2.

All random factors and coordinate indices can be drawn in advance. The
query choice does not depend on a, y, or any other response. Arithmetic
and memory can be O(Nm) or larger, but these are not charged in RA-11.

\subsubsection{7.3 A regime in which the diagonal branch improves the
preceding
envelope}\label{a-regime-in-which-the-diagonal-branch-improves-the-preceding-envelope}

Take n=2 and epsilon\_q=(9/10)\^{}q. Up to universal constants and
polynomial factors, the previous displayed upper envelope in this
conversation gives

\begin{verbatim}
[(4/3)/(9/10)^2]^q = (400/243)^q,
\end{verbatim}

whose base is approximately 1.64609; this branch is below the exact
2\^{}q bound. The earlier linear-accuracy hierarchy has base greater
than 2 in this binary case and does not improve that comparison.

The new diagonal estimator gives

\begin{verbatim}
O[(sqrt(2)/(9/10))^q]
  = O[(10sqrt(2)/9)^q],
\end{verbatim}

with base approximately 1.57135. The ratio of the new base to the old is
27sqrt(2)/40\textless1 (squaring reduces this to 1458\textless1600).
This is an asymptotic comparison of proved query counts, not a claim
that an exponential-size matrix experiment was run.

\subsection{8. What is and is not
resolved}\label{what-is-and-is-not-resolved}

The new upper envelope is the minimum of the previously reported bounds,
the diagonal bound of Theorem 2.2, and the infimum of Theorem 2.1 over p
in (1,2{]}. The finite-bit version is available with at most a factor
two in sample count.

At n=2 and epsilon=1/4, the previously retained blocking lower bound has
scale

\begin{verbatim}
Omega(sqrt(q/log q)),
\end{verbatim}

whereas the new upper bound has scale

\begin{verbatim}
exp[(log 2 - 1/2)q + O(sqrt q)] times a polynomial in q.
\end{verbatim}

This improves the previous exponential upper rate but leaves an
exponential-versus-polynomial gap. Theorem 5.1 does not close it: its
lower bound is only for an empirical-average estimator, and its hard
rank-one example is easy for the full-vector oracle.

Accordingly, this package does not determine Q to universal constant
factors jointly in n,q,epsilon and is not a full solution of RA-11. The
missing statement is a matching unrestricted lower bound, a
substantially stronger unrestricted upper bound, or another
characterization closing the gap. No repository status has been changed.

\subsection{9. Reproducibility and
limitations}\label{reproducibility-and-limitations}

\texttt{code/exact\_checks.py} checks exact rational diagonal
identities, coordinate correction identities, and Gaussian-rational real
interpolation on finite examples. These calculations catch
algebraic/implementation mistakes but are not formal proofs of the
general theorems.

\texttt{code/numerical\_checks.py} runs finite diagonal-estimator
experiments and scalar Gamma-law diagnostics for the empirical-average
estimator. The latter do not use a giant matrix and must not be
described as full-oracle lower-bound experiments. Query hashes check
nonadaptivity of the diagonal implementation at fixed random seeds. The
numerical script records its actual parameter choices, query counts, and
failures.

Floating-point linear algebra, interpolation, and random variate
routines are not exact arithmetic. The finite-support fractional-moment
theorem is a constructive existence argument; a fully certified
discretization implementation is not supplied. No floating-point
stability guarantee, efficient runtime theorem, independent mathematical
review, or proof-assistant certification is claimed.

\subsection{References and provenance}\label{references-and-provenance}

{[}R1{]} RA-11 repository statement, OpenProblemsInNLA,
randomized-and-low-rank-approximation/RA-11/README.md. Public source:
\url{https://raw.githubusercontent.com/ajt60gaibb/OpenProblemsInNLA/main/randomized-and-low-rank-approximation/RA-11/README.md}
. See the local source snapshot and retrieval metadata when available.

{[}R2{]} Prior conversation manuscript,
\texttt{ra11\_partial\_results.pdf}, round 1. Source of the reported
blocking lower bound and the original real-query interpolation
construction in this research sequence.

{[}R3{]} Prior conversation manuscript,
\texttt{ra11\_linear\_accuracy.pdf}, round 2. Reported
fixed-tensor-order optimality for arbitrary PSD inputs and the first
general linear-accuracy hierarchy.

{[}R4{]} Prior conversation manuscript,
\texttt{ra11\_superset\_controls.pdf}, round 3. Reported
superset-specific controls and the improved linear-accuracy upper bound.
This package does not assert an independent re-verification of all
inherited proofs.

The new arguments are written out above. No claim is made that their
ingredients, such as importance-sampling entropy rates or diagonal
control variates, are new to all literature.

\end{document}
